\chapter*{Convenciones del manual}
\addcontentsline{toc}{chapter}{Convenciones del manual}

\begin{nota}
Las capturas finales usan datos ficticios y marcadores numerados discretos. Las explicaciones largas se ubican en el texto o en el pie de figura, no dentro de la imagen.
\end{nota}

\begin{tabularx}{\linewidth}{>{\bfseries}p{3.6cm}X}
\toprule
Elemento & Uso en el manual\\
\midrule
Paso numerado & Indica una acción que debe ejecutarse en orden.\\
Nota & Aclara una recomendación o condición útil que no bloquea el procedimiento.\\
Advertencia & Señala riesgos de datos personales, envío de información o confusión operativa.\\
Operación sensible & Marca una acción que puede cambiar resultados, asistencia o información importante.\\
Resultado esperado & Describe qué debe observar el usuario al terminar un procedimiento.\\
Campo obligatorio & Se menciona cuando el formulario no debe enviarse vacío o con información incompleta.\\
Botón & Se escribe en negrita, por ejemplo \botoninterfaz{Enviar inscripción}.\\
Sección de pantalla & Se nombra como aparece en la interfaz, por ejemplo \campointerfaz{Reglamento}.\\
Marcador de captura & Un círculo numerado identifica zonas de una imagen; la explicación aparece bajo la figura o en el texto.\\
\bottomrule
\end{tabularx}

\section*{Lectura de capturas}

Las figuras muestran el estado esperado de una pantalla. Si una figura tiene marcadores, cada número se explica en el procedimiento correspondiente. Cuando una misma captura sirve para varios procedimientos, se referencia la figura ya presentada en lugar de repetirla.

\section*{Rutas y ambientes}

Las rutas se escriben como ejemplos relativos, por ejemplo \ruta{/registro/}. La URL de producción queda como \textbf{PENDIENTE DE DEFINICIÓN PARA PRODUCCIÓN}.
